\section{扩散模型为何能建模复杂分布}

\subsection{单步高斯 vs. 多步高斯}

讲者在课堂上提出了一个深刻的问题：前向过程的转移是高斯的，反向过程的转移也被建模为高斯的，那么整个似然分布 $p_\theta(x_0|x_T)$ 是否也是高斯分布？如果是，那扩散模型岂不是和 VAE 一样受限？

答案是\textbf{否}。

\begin{importantbox}{非线性带来复杂性}
虽然每一步反向转移 $p_\theta(x_{t-1}|x_t)$ 都是高斯分布，但整体似然分布 $p_\theta(x_0|x_T) = \int p_\theta(x_0|x_1) p_\theta(x_1|x_2) \cdots p_\theta(x_{T-1}|x_T) \, dx_{1:T-1}$ \textbf{不是}高斯分布。关键在于：每一步高斯的\textbf{均值}是由神经网络（非线性函数）预测的。

与前向过程不同——前向过程中均值是 $x_{t-1}$ 的线性函数（$\sqrt{1-\beta_t}\, x_{t-1}$），因此前向过程的多步组合仍然是高斯的。而反向过程中均值是 $x_t$ 的非线性函数（通过神经网络），这使得多步组合可以表达任意复杂的分布。
\end{importantbox}

\begin{warningbox}{线性前向 vs. 非线性反向}
这是理解扩散模型表达能力的关键：
\begin{itemize}
    \item \textbf{前向过程}：均值是输入的\textbf{线性}函数 $\rightarrow$ 多步组合仍为高斯 $\rightarrow$ 可推导闭式跳跃分布
    \item \textbf{反向过程}：均值是输入的\textbf{非线性}函数（神经网络）$\rightarrow$ 多步组合\textbf{不是}高斯 $\rightarrow$ 似然分布可以任意复杂
\end{itemize}
扩散模型巧妙地在``训练每一步简单''和``整体表达能力强''之间找到了平衡。
\end{warningbox}

这也解释了为什么扩散模型需要多步：步数越多，非线性变换的层次越深，表达能力越强。这类似于深度神经网络比浅层网络更有表达力的道理。

\subsection{本章小结}

扩散模型的每一步转移虽然是高斯的，但由于反向过程使用非线性神经网络预测均值，多步组合的整体似然分布可以建模任意复杂的分布。这是扩散模型优于简单 VAE 的根本原因。

%% ========== 第九节：噪声预测与 Score Matching 的联系 ==========

